\PassOptionsToPackage{usenames,dvipsnames}{xcolor}
\documentclass[a4paper,12pt]{letter}
\usepackage[utf8]{inputenc}
\usepackage[T1]{fontenc}
\usepackage[english,slovene]{babel}
\usepackage[tikz]{bclogo}
\usepackage{qrcode,marvosym}
\usepackage{pgf,tikz,pgfplots}
\pgfplotsset{compat=1.14}
\usepackage{mathrsfs}
\usetikzlibrary{arrows}
\usepackage{graphicx}
\newcommand{\ds}{\displaystyle}
\usepackage{fancybox}
\usepackage{amsfonts}
\usepackage{amsmath}
\usepackage{wallpaper}
\usepackage{hyperref}
\usepackage{array}
\newcolumntype{M}[1]{>{\centering\arraybackslash}m{#1}}
\newcolumntype{N}{@{}m{0pt}@{}}
\usepackage{fancyhdr,enumerate}
\newcommand{\ora}{\overrightarrow}
\renewcommand{\headrulewidth}{3pt}
\renewcommand{\headrule}{\hbox to\headwidth{%
  \color{gray}\leaders\hrule height \headrulewidth\hfill}}
\renewcommand{\footrulewidth}{2pt}
\renewcommand{\footrule}{\hbox to\headwidth{%
  \color{brown}\leaders\hrule height \footrulewidth\hfill}}
\usepackage{gensymb}
\usepackage{amssymb}
  \textwidth 20 true cm
  \oddsidemargin -2 true cm
  \evensidemargin -2 true cm
  \topmargin -3.5 true cm
  \textheight 27.5 true cm
  \linespread{1.2}
\renewcommand{\headrulewidth}{0.3pt}
\renewcommand{\footrulewidth}{1.9pt}
\definecolor{qqwuqq}{rgb}{1,0,0}
\definecolor{qqqqff}{rgb}{0,0,1}
\definecolor{qqffqq}{rgb}{0,1,0}
\definecolor{ffqqqq}{rgb}{1,0,0}
\definecolor{zzuxux}{rgb}{0,0,0}
\usepackage{mdframed}
\usepackage{multicol}
\definecolor{bgblue}{RGB}{0,0,253}
\usepackage{eurosym}
%%%%%%%%%%%%%%%%%%%%%%%%%%%%%%%%%%%%%%%%%%%%%%%%%%%%%%%%
\newcounter{tocke}
\setcounter{tocke}{0}
\usepackage{tikz-3dplot}
%%%%%%%%%%%%%%%%%%%%%%%%%%%%%%%%%%%%%%%%%%%%%%%%%%%%%%%%%%
\mdfdefinestyle{theoremstyle}{%
linecolor=brown!40,linewidth=1pt,%
frametitlerule=true,%
frametitlebackgroundcolor=brown!50,
innertopmargin=\topskip,
}
\mdfdefinestyle{theoremstyle1}{%
linecolor=brown,middlelinewidth=1pt,%
frametitlerule=true,%
apptotikzsetting={\tikzset{mdfframetitlebackground/.append style={%
shade,left color=black!40, right color=gray!10},
mdfbackground/.append style={
top color=white!100!white,
bottom color=black!5
},
}},
frametitlerulecolor=gray,
frametitlerulewidth=2pt,
innertopmargin=\topskip,
innerbottommargin=\topskip,
}
\mdtheorem[style=theoremstyle1]{naloga}{Naloga}
\mdtheorem[style=theoremstyle]{resitev}{Analiza Naloge}
\mdtheorem[style=theoremstyle]{kriterij}{\v{S}tevilo dose\v{z}enih to\v{c}k na testu}
%%%%%%%%%%%%%%%%%%%%%%%%%%%%%%%%%%%%%%%%%%%%%%%%%%%%%%%%%%
\usepackage[table]{xcolor}
\usepackage[printwatermark]{xwatermark}
\newwatermark[allpages,color=brown!2,angle=45,scale=5,
xpos=-5,ypos=0]{matej.info}
%%%%%%%%%%%%%%%%%%%%%%%%%%%%%%%%%%%%%%%%%%%%%%%%%%%%%%%%%%

\begin{document}
\pagestyle{fancy}
\renewcommand\headrulewidth{0pt}
\lhead{}\chead{}\rhead{}
\rfoot{\vspace*{-2\baselineskip}}
\renewcommand{\vec}{\overrightarrow}
%%%%%%%%%%%%%%%%%%%% podatki testa -- tu spreminjaj
\def\programletnik{Gim - 4}
\def\naslov{Test 1.4.P - }
\def\sklop{Geometrijska telesa}
%%%%%%%%%%%%%%%%%%%% tocke po nalogah
\def\tockeA{3+4}
\def\tockeB{4+4}
\def\tockeC{3+3+4}
\def\tockeD{3+3+4}
%%%%%%%%%%%%%%%%%%%%
\begin{minipage}{20cm}
\begin{bclogo}[couleur=brown!30,
  arrondi=0,ombre=true, couleurOmbre=gray!50,
logo =\bcplume,
  noborder=true]
{\textrm{{\Large{\naslov \sklop} \Huge{}
\large{} \hfill $\programletnik$}}}
\end{bclogo}
\end{minipage}
\begin{minipage}{20cm}
\begin{bclogo}[couleur=brown!30,
  arrondi=0,ombre=true, couleurOmbre=gray!50,
logo =\bcplume,
  noborder=true]
{\textrm{{\raisebox{-0.3cm}{\large{Ime in Priimek: \rule{10cm}{0.1mm}} \Huge{}
\large{} \hfill $ $}}}}
\end{bclogo}
\end{minipage}
%%%%%%%%%%%%%%%%%%%%%%%%%%%%%%%%%%%%%%%%%%%%%%%%%

%%%%%%%%%%%%%%%%%%%%%%%%%%%%%%%%%%%%%%%%%%%%%%%%%
\begin{naloga}[\hfill $\tockeA$] \addtocounter{tocke}{\numexpr\tockeA\relax}
\v{S}tiristrani pokon\v{c}ni prizmi je za osnovno ploskev pravokotni trapez z vzporednima stranicama $9$ cm in $6$ cm, pravokotnim krakom (vi\v{s}ino trapeza) $4$ cm in po\v{s}evnim krakom $5$ cm. Rob (vi\v{s}ina) prizme meri $9$ cm.
\begin{enumerate}[a)]
\item Izra\v{c}unaj plo\v{s}\v{c}ino osnovne ploskve.
\item Izra\v{c}unaj prostornino in povr\v{s}ino prizme.
\end{enumerate}
\end{naloga}\newpage

%%%%%%%%%%%%%%%%%%%%%%%%%%%%%%%%%%%%%%%%%%%%%%%%%
\begin{naloga}[\hfill $\tockeB$] \addtocounter{tocke}{\numexpr\tockeB\relax}
Pravilna \v{s}esterokotna piramida ima osnovni rob $a=6$ cm in vi\v{s}ino $v=4\sqrt{3}$ cm.
\begin{enumerate}[a)]
\item Izra\v{c}unaj prostornino piramide.
\item Izra\v{c}unaj povr\v{s}ino piramide.
\end{enumerate}
\end{naloga}\newpage

%%%%%%%%%%%%%%%%%%%%%%%%%%%%%%%%%%%%%%%%%%%%%%%%%
\begin{naloga}[\hfill $\tockeC$] \addtocounter{tocke}{\numexpr\tockeC\relax}
Sto\v{z}ec ima polmer osnovne ploskve $r=12$ cm in vi\v{s}ino $v=16$ cm.
\begin{enumerate}[a)]
\item Izra\v{c}unaj stranico sto\v{z}ca in njegovo povr\v{s}ino.
\item Izra\v{c}unaj prostornino sto\v{z}ca.
\item Izra\v{c}unaj kot med stranico in vi\v{s}ino (osjo) sto\v{z}ca, do kotne minute natan\v{c}no.
\end{enumerate}
\end{naloga}\newpage

%%%%%%%%%%%%%%%%%%%%%%%%%%%%%%%%%%%%%%%%%%%%%%%%%
\begin{naloga}[\hfill $\tockeD$] \addtocounter{tocke}{\numexpr\tockeD\relax}
Krogla in valj imata enako prostornino. Polmer krogle je $R=6$ cm, polmer valja pa je enak polmeru krogle.
\begin{enumerate}[a)]
\item Izra\v{c}unaj prostornino krogle.
\item Izra\v{c}unaj vi\v{s}ino valja.
\item V kak\v{s}nem razmerju sta povr\v{s}ini krogle in valja (pla\v{s}\v{c} in obe osnovni ploskvi)?
\end{enumerate}
\end{naloga}

\newpage
\vfill
\begin{kriterij*} \begin{center} \renewcommand{\arraystretch}{2.5}
\begin{tabular}{|>{\columncolor{gray!20}}c||c|c|c|c|c||c|c|c|} \hline Odstotki &$[0,45)$ & $[45,60)$ & $[60,75)$ & $[75,90)$ & $[90,100]$& \v{S}t. mo\v{z}nih to\v{c}k& Dose\v{z}ene to\v{c}ke & OCENA \\ \hline Ocena & 1 & 2 & 3 & 4 & 5 &\cellcolor{gray!20} $\thetocke$&\cellcolor{gray!20} \phantom{ocena} & \cellcolor{gray!20} \phantom{a} \\ \hline \end{tabular} \end{center} \vspace{0.3cm}  \end{kriterij*}

\newpage
\textsc{Re\v{s}itve.}

\begin{enumerate}
\item a) plo\v{s}\v{c}ina osnovne ploskve $=30\ \text{cm}^2$.\quad b) $V=270\ \text{cm}^3$, $S=276\ \text{cm}^2$.
\item a) $V=216\ \text{cm}^3$.\quad b) $S=144\sqrt3\ \text{cm}^2\approx249{,}42\ \text{cm}^2$.
\item a) $s=20$ cm, $S=384\pi\ \text{cm}^2\approx1206{,}37\ \text{cm}^2$.\quad b) $V=768\pi\ \text{cm}^3\approx2412{,}74\ \text{cm}^3$.\quad c) $\varphi\approx36^{\circ}52'$.
\item a) $V_{\text{krogla}}=288\pi\ \text{cm}^3\approx904{,}78\ \text{cm}^3$.\quad b) $v_{\text{valj}}=8$ cm.\quad c) $S_{\text{krogla}}:S_{\text{valj}}=6:7$ ($S_{\text{krogla}}=144\pi$, $S_{\text{valj}}=168\pi$).
\end{enumerate}
\end{document}
